\documentclass[10pt,a4paper]{amsart}
\usepackage[margin=19mm]{geometry}
\usepackage{amsmath,amssymb,mathtools}
\usepackage[scr=rsfs,cal=euler]{mathalfa}
\usepackage{xfrac}
\usepackage[unicode,hidelinks,bookmarks=false]{hyperref}
\newcommand{\ie}{i.e.\ }
\newcommand{\cf}{cf.\ }
\newcommand{\resp}{respectively\ }
\newcommand{\Subsection}{\subsection}
\providecommand{\nobreakdash}{\nobreak\hbox{-}}
\setlength{\emergencystretch}{2em}
\multlinegap0pt
\usepackage{tikz}
\usepackage{amssymb}
\usepackage{float}
\usepackage{xcolor}
\usepackage[shortlabels]{enumitem}
\usepackage{enumerate}
\newtheorem{theorem}{Theorem}
\def\Des{\mathsf{Des}}
\def\cyc{\mathsf{cyc}}
\def\rmin{\mathsf{rmin}}
\title{On two conjectures concerning special kinds of descents on permutations}
\author{Taifeng Ding, Sherry H.F. Yan}
\date{Source: arXiv:2608.15521v1 (CC BY 4.0)}
\begin{document}
\maketitle

\noindent\emph{Source and licence.} On two conjectures concerning special kinds of descents on permutations. Version arXiv:2608.15521v1 (CC BY 4.0), distributed by arXiv under the Creative Commons Attribution 4.0 International licence, http://creativecommons.org/licenses/by/4.0/, as recorded in the arXiv licence metadata for this exact version and shown on the abstract page https://arxiv.org/abs/2608.15521v1. Full licence notice: \texttt{licenses/arxiv-2608.15521-CC-BY-4.0.txt}.\par\smallskip

\noindent\textit{Editorial scope.} Counted unit: the article's theorem thm-Des\_2 (source0.tex lines 364--369). The article's complete original proof of it is delivered with it (source0.tex lines 370--417, one \texttt{proof} environment of 48 lines). No mathematics is written, repaired or re-derived: the statement and the proof are the article's own lines, in the article's own notation, and no retained line is substituted or re-flowed.  No further source-local context is needed: every cross-reference of every retained line resolves inside the two delivered spans.  Nothing was omitted from any retained span, so the package carries no omission record.  No retained line contains a citation, so the wrapper carries no bibliography and no bibliography entry.  No retained line contains a figure environment, an image-inclusion command or drawing code, so no image is delivered and the package declares no asset dependency.  Editorial formatting only, in addition: an AMS \texttt{amsart} preamble that restates the article's own declarations and macros used by the retained lines, namely the environments \texttt{theorem} on separate counters, together with the standard packages those lines need.  The retained environments print in this document as Theorem 1 in order of appearance, whereas the article prints the same items with the numbering of its own full document.  The complete native source of the article as distributed in the arXiv e-print for this exact version is bundled unchanged as \texttt{sources/207756/source0.tex}.
     \begin{theorem}\label{thm-Des_2}
     The map $\delta:=\chi\circ\theta$ induces a bijection from $\mathfrak{S}_n$ onto itself such that for any permutation $\pi\in \mathfrak{S}_n$, we have
     	\begin{equation}\label{eq-Des_2}
     	(\Des_2, \rmin)\pi= (\Des_2, \cyc)\delta(\pi).
     		\end{equation}
     	\end{theorem}

     \begin{proof}
    Since both   $\theta$ and $\chi$ are bijections, the map $\delta$   is a bijection from   $\mathfrak{S}_n$ onto itself.  Now we proceed to  show that \eqref{eq-Des_2} holds for any $\pi\in \mathfrak{S}_n$. Let $\pi=\pi(1)\pi(2)\cdots\pi(n)$, $\theta(\pi)=s=(a_1,a_2,\ldots,a_n)$ and $\chi(s)=\sigma=\sigma(1)\sigma(2)\cdots\sigma(n)$. When we apply the map $\theta^{-1}$ to $s=(a_1,a_2,\ldots,a_n)$, we generate a sequence of permutations $\pi^{(1)},\pi^{(2)},\ldots,\pi^{(n)}$ such that $\pi^{(n)}=\pi$. Recall that $\pi^{(1)}=1$, and for all $i>1$, the permutation $\pi^{(i)}$ is obtained from $\pi^{(i-1)}$ by inserting the letter $i$ immediately before the $a_i$-th letter of $\pi^{(i-1)}$ if $a_i<i$, and putting the letter $i$ at the end of $\pi^{(i-1)}$ otherwise.
     
     Similarly, when applying the map $\chi$ to $s$, we also obtain a sequence of permutations $\sigma^{(1)},\ldots,\sigma^{(n)}$, where $\sigma^{(1)}=1$ and $\sigma^{(n)}=\sigma$. For  each  $i>1$,   $\sigma^{(i)}$ is determined by the following rules.
     \begin{itemize}
     	\item If $a_i = i$, define   $\sigma^{(i)}$ to be the permutation  whose cycle notaion is obtained from that  of $\sigma^{(i-1)}$ by inserting the cycle $(i)$.
     	\item Otherwise, define   $\sigma^{(i)}$ to be the permutation  whose cycle notaion is obtained from that of   $\sigma^{(i-1)}$ by  inserting  the entry $i$ immediately after the entry $a_i$. 
     \end{itemize}
     
     In what follows, we aim to show that
     \[
     (\Des_2, \rmin)\pi^{(i)}= (\Des_2, \cyc)\sigma^{(i)}
     \]
     holds for all $1\leq i\leq n$. We proceed by induction on $i$. Clearly, the assertion holds for $i=1$ since $\pi^{(1)}=\sigma^{(1)}=1$. Now assume that $i>1$, and suppose that
     \[
     (\Des_2, \rmin)\pi^{(i-1)}= (\Des_2, \cyc)\sigma^{(i-1)}.
     \]
     
     By the construction of $\pi^{(i)}$, it is easily seen that
     \[
     \Des_2(\pi^{(i)})= \{k\mid  k\in \Des_2(\pi^{(i-1)}) \;\; \mbox{and}\;\; k<a_i\}  \cup\{a_i\},\quad \rmin(\pi^{(i)})=\rmin(\pi^{(i-1)})
     \]
     when $a_i<i$, and
     \[
     \Des_2(\pi^{(i)})= \Des_2(\pi^{(i-1)}),\quad \rmin(\pi^{(i)})=\rmin(\pi^{(i-1)})+1
     \]
     otherwise.
     
     Similarly, by the construction of $\sigma^{(i)}$, one can easily check that if $a_i=i$, then $\sigma^{(i)}(j)=\sigma^{(i-1)}(j)$ for all $j<i$ and $\sigma^{(i)}(i)=i$. Moreover, if $a_i<i$, then $\sigma^{(i)}(j)=\sigma^{(i-1)}(j)$ for all $j<i$ and $j\neq a_i$, $\sigma^{(i)}(a_i)=i$ and $\sigma^{(i)}(i)=\sigma^{(i-1)}(a_i)$. Therefore, we deduce that
     \[
     \Des_2(\sigma^{(i)})= \{k\mid  k\in \Des_2(\sigma^{(i-1)})\;\; \mbox{and}\;\; k<a_i \}\cup\{a_i\},\quad \cyc(\sigma^{(i)})=\cyc(\sigma^{(i-1)})
     \]
     when $a_i<i$, and
     \[
     \Des_2(\sigma^{(i)})= \Des_2(\sigma^{(i-1)}),\quad \cyc(\sigma^{(i)})=\cyc(\sigma^{(i-1)})+1
     \]
     otherwise.
     
     Then, by the induction hypothesis, we conclude that
     \[
     \Des_2(\pi^{(i)})=\Des_2(\sigma^{(i)})\quad \mbox{and}\quad \rmin(\pi^{(i)})=\cyc(\sigma^{(i)})
     \]
     for all $1\leq i\leq n$. Hence, we deduce that
     \[
     (\Des_2, \rmin)\pi=  (\Des_2, \rmin)\pi^{(n)}= (\Des_2, \cyc)\sigma^{(n)}=(\Des_2, \cyc)\sigma
     \]
     as desired, completing the proof. 
     	\end{proof}

\end{document}
